% =====================================================================
%  CHAPTER 1: 绪论（对应大纲第一章：绪论）
%  参考PPT：第一章 绪言（郝加虎，安徽医科大学）
% =====================================================================
\chapter{绪论}
\setcounter{chapter}{1}
\chapterintro{
\textbf{掌握：}儿童少年生长发育的规律及其卫生学意义。\newline
\textbf{熟悉：}课程的性质与任务、研究方法和应用体系。\newline
\textbf{了解：}学科发展历程与趋势。
}

% ----- 1.1 儿童少年卫生学概述 -----
\section{儿童少年卫生学概述}

\subsection{研究对象}

\begin{defbox}
\textbf{儿童少年卫生学（child and adolescent health）：}研究维护和促进儿童少年健康的一门学科。研究儿童少年身心发育随年龄变化的特征，分析影响生长发育的遗传和环境因素，阐明儿童少年机体与学习及生活环境间的相互关系，制定相应卫生要求和措施，预防疾病、增强体质，促进身心健康发育，并为成年期健康奠定良好基础，从而达到提高生命质量的目的。

\textbf{研究对象：}0-24岁儿童少年；重点人群：中小学生群体。
\end{defbox}

\subsection{儿童少年年龄范围的界定}

\begin{keybox}
\textbf{儿童少年年龄分期：}

\begin{table}[H]
\centering
\caption{儿童少年年龄分期}
\begin{tabularx}{\textwidth}{l>{\raggedright\arraybackslash}p{3.5cm}X}
\hline
\textbf{分期} & \textbf{英文名称} & \textbf{年龄范围} \\
\hline
\imp{胎儿期} & fetal period & 从受精卵形成开始到孕40周胎儿娩出 \\
\imp{婴儿期} & infant period & 出生后1年 \\
\imp{新生儿期} & neonatal period & 出生后28天（属于婴儿期） \\
\imp{幼儿期} & toddle period & 出生后第2-3年 \\
\imp{学龄前期} & preschool period & 3-6岁 \\
\imp{学龄期} & school-age period & 6-11/12岁，相当于小学阶段 \\
\imp{青春期} & adolescence & 从青春发动开始到生长基本成熟（10-19岁） \\
\imp{青年期} & youth & 15-24岁 \\
\hline
\end{tabularx}
\end{table}
\end{keybox}

\subsection{儿童少年生物和社会特征}

\begin{keybox}
\textbf{儿童少年三大生物和社会特征：}

\begin{enumerate}[leftmargin=*]
    \item \imp{处于生长发育过程中}——生长发育是儿童少年的基本特征之一，也是社会发展的一面镜子。了解儿童少年生长发育规律，并深入研究其影响因素，是制定儿童少年保健工作的相应政策和行动纲领的重要前提。

    \item \imp{接受学校教育的塑造}——接受教育是儿童少年的一个重要社会特征，学校集体生活构成其重要的生活内容。创造良好的学校物质环境和社会支持环境，建立健康促进学校，开展综合性学校卫生服务，是促进其身心健康和生长发育潜能的重要条件。

    \item \imp{具有社会脆弱性和健康易损性}——在全人口中，儿童少年由于其身心以及各种能力发育的不完善，始终是一个弱势的群体。儿童神经系统发育不成熟，与成人相比，更容易受到有害生物和环境因素影响，从而造成发育和行为异常。在不同的社会生态环境层面——个体、人际间、社区、结构和政策层面，以及在不同的发展阶段，儿童少年的心理健康受到多个动态因素的影响。
\end{enumerate}
\end{keybox}

\subsection{儿童少年生存与发展权利}

\begin{defbox}
\textbf{儿童少年生存与发展权利：}

\begin{enumerate}[leftmargin=*]
    \item \imp{生存权}：每一个儿童都应该享有的对自己生命权以及维持基本健康生活需求的生活条件保障权。

    \item \imp{发展权}：保障儿童成长过程中的各种需要得到满足，包括儿童有接受一切形式的教育（正规的和非正规的教育）的权利，能够给予儿童的身体、心理、精神、道德与社交发展的相应生活水平。发展权包括参与权、受保护权。
\end{enumerate}
\end{defbox}

% ----- 1.2 课程的性质与任务 -----
\section{课程的性质与任务}

\begin{defbox}
\textbf{课程性质：}儿童少年卫生学是预防医学专业学生的一门必修考试课，是研究维护和促进儿童少年健康的重要专业课程。该课程以儿童少年的身心发育为主线，综合应用预防医学、临床医学、心理学、教育学等多学科知识，具有鲜明的应用性和实践性。

\textbf{课程任务：}研究儿童少年生长发育的特点和规律；分析影响儿童少年生长发育和健康的遗传因素与环境因素；阐明儿童少年机体与学习、生活环境之间的相互关系；制定相应的卫生要求和干预措施；预防疾病、增强体质；促进儿童少年身心健康发展；为成年期健康奠定良好基础；达到提高生命质量和健康水平的目的。
\end{defbox}

% ----- 1.3 生长发育规律 -----
\section{儿童少年生长发育的规律及其卫生学意义}

\begin{keybox}
\textbf{儿童少年生长发育的规律：}

生长发育是儿童少年的基本特征之一，也是社会发展的一面镜子。儿童少年身心发育随年龄变化的特征，受遗传和环境因素共同影响。生长发育是累积和连续的过程，个体从"胎儿—儿童—少年—青壮年—老年"构成完整的生命历程，每一个阶段健康既是前一阶段的延续，也是下一阶段的基础。

\textbf{卫生学意义：}

\begin{enumerate}[leftmargin=*]
    \item 了解儿童少年生长发育规律，并深入研究其影响因素，是制定儿童少年保健工作的相应政策和行动纲领的重要前提。
    \item 儿童少年发展潜力的发挥必须通过整个生命历程中的支持性环境来实现。"支持性环境"是能够激发儿童少年最佳发展潜力的系统、力量和因素，是由贯穿整个生命历程的卫生服务、政策、计划和社区共同参与。
    \item 生长发育规律为制定相应卫生要求和措施、预防疾病、增强体质、促进身心健康发育提供科学依据，并为成年期健康奠定良好基础，从而达到提高生命质量的目的。
\end{enumerate}
\end{keybox}

% ----- 1.4 应用体系与学科发展 -----
\section{应用体系与学科发展}

\subsection{应用体系}

\begin{defbox}
\textbf{儿童少年卫生学的应用体系：}

本学科的应用体系覆盖了从个体到群体的多层次健康服务，主要包括以下组成部分：

\begin{enumerate}[leftmargin=*]
    \item \imp{学校卫生服务体系}——以学校为基础，为在校儿童少年提供综合性卫生服务，包括健康检查、疾病筛查、心理咨询、营养指导等，是学科最重要的实践领域。

    \item \imp{健康监测与健康监督}——通过建立常态化的儿童少年健康监测体系，系统收集生长发育、疾病发生、健康行为等数据，为卫生决策提供科学依据。包括学生体质健康监测、常见病监测、健康危险行为监测等。

    \item \imp{健康教育与健康促进}——以学校为主要场所，通过课程教育、环境创设、政策支持等综合策略，提高儿童少年的健康素养，培养健康行为，创建有利于健康发展的支持性环境。

    \item \imp{疾病预防与控制}——针对儿童少年常见病、传染病、伤害、心理行为问题等，实施三级预防策略，构建学校传染病防控体系和伤害防控体系。

    \item \imp{卫生政策制定与评估}——参与制定国家及地方的儿童少年卫生政策、法规和标准，推动学校卫生工作的制度化和规范化，并对政策实施效果进行评估。
\end{enumerate}

\textbf{应用体系的整合：}本学科的应用依托于\textbf{三大系统的协同}：\imp{妇幼保健体系}（覆盖0-6岁儿童）、\imp{学校卫生体系}（覆盖6-18岁学龄儿童少年）、\imp{社区卫生服务体系}（提供全生命周期的基层卫生保健），三者相互衔接、协调配合，共同构成儿童少年健康的保障网络。
\end{defbox}

\subsection{学科发展历程与趋势}

\begin{tipbox}
\textbf{学科发展历程与趋势：}

儿童少年卫生学已从传统的学校卫生发展为现代综合性儿童少年健康科学。学科发展经历了以下关键阶段：

\begin{enumerate}[leftmargin=*]
    \item \imp{学科独立与奠基阶段}——20世纪初，随着学校卫生概念的引入，儿童少年卫生学逐渐从公共卫生学科中分化出来，成为独立的学科领域。我国于20世纪50年代正式确立\"学校卫生学\"的学科地位，后更名为\"儿童少年卫生学\"。

    \item \imp{多学科融合阶段}——将预防医学、临床医学、心理学、教育学、社会学等多学科理论和方法融入儿童少年健康研究与实践，学科内涵不断丰富。

    \item \imp{生命历程视角的引入}——提出\"生命历程观点\"（life course perspective），强调从胎儿期到成年期的全程健康管理。引入\textbf{DOHaD理论}（Developmental Origins of Health and Disease，健康与疾病的发育起源学说），阐明生命早期（孕期和出生后早期）环境暴露对终身健康的关键编程效应。

    \item \imp{社会生态模型的应用}——建立社会生态模型（social-ecological model）分析框架，从个体、人际、社区、政策等多层面系统分析影响儿童少年健康的因素，推动多层次、综合性干预策略的制定。

    \item \imp{健康促进学校框架的建立}——世界卫生组织（WHO）倡导的\imp{健康促进学校}（health-promoting schools）框架，将学校建设为促进儿童少年健康发展的综合性平台，强调课程教学、学校环境、卫生服务和社区参与的一体化整合。

    \item \imp{循证干预的推广}——强调基于科学证据的干预措施和卫生决策，运用流行病学和统计学方法评估干预效果，推动学科从经验模式向循证模式转变。

    \item \imp{健康中国战略的推动}——《\"健康中国2030\"规划纲要》将儿童少年健康列为重点领域，明确提出加强学校卫生工作、促进儿童早期发展、保障青少年身心健康的战略目标，为学科发展提供了前所未有的政策支持和实践平台。
\end{enumerate}

\textbf{学科发展趋势：}未来将更加注重生命历程的纵向整合、多学科交叉融合、大数据与精准健康管理、以及全球儿童健康治理的参与。
\end{tipbox}

% ----- 1.5 儿童少年卫生学学科体系 -----
\section{儿童少年卫生学学科体系}

\subsection{研究内容}

\begin{defbox}
\textbf{儿童少年卫生学的四大研究内容：}

\begin{enumerate}[leftmargin=*]
    \item \imp{儿童少年生长发育规律}——生长发育作为人类最为复杂的生物学现象，是本学科研究的永恒主题，也是学科首先要研究和认识的理论问题。包括身体发育和心理发育两个方面。

    \item \imp{儿童少年健康状况及其决定因素}——以保护、促进、增强儿童少年身心健康为宗旨，提出相应的卫生要求和适宜的卫生措施，维护儿童少年的健康，减少疾病，预防死亡，提高生存质量。

    \item \imp{儿童少年卫生服务}——包括：生长发育和健康监测；健康教育和健康促进；常见病防治与健康管理；心理卫生服务；学校营养服务；校园安全管理与伤害及暴力预防；体育与体力活动促进。

    \item \imp{儿童少年卫生学的技术和方法}——人体测量、健康相关体能测试、行为评定等。
\end{enumerate}
\end{defbox}

\subsection{研究方法}

\begin{defbox}
\textbf{儿童少年卫生学的主要研究方法：}

\begin{enumerate}[leftmargin=*]
    \item \imp{人体测量学方法（anthropometry）}——通过规范化的测量工具和测量方法，对儿童少年的身体形态指标进行定量评估，包括身高、体重、坐高、胸围、头围、皮褶厚度等基础指标的测量，以及体型评定、身体成分分析等。人体测量是本学科最经典、最基础的研究手段，为生长发育评价提供核心数据。

    \item \imp{生理功能检测方法}——通过检测儿童少年的生理功能指标，评价其机能发育水平和健康状况，主要包括：心肺功能测试（肺活量、最大摄氧量、心率变异性等）、内分泌功能检测（激素水平测定）、神经生理功能检查（脑电图、诱发电位等）、体能测试（肌力、耐力、柔韧性、速度、协调性等）。

    \item \imp{心理行为评定方法}——运用标准化心理测量工具评估儿童少年的心理发育水平和行为特征，主要包括：
        \begin{itemize}
            \item \textbf{智力测验}：如韦氏儿童智力量表（WISC）、瑞文推理测验等，评估一般认知能力。
            \item \textbf{行为评定量表}：如Achenbach儿童行为量表（CBCL）、Conners儿童行为问卷、长处与困难问卷（SDQ）等，用于筛查和评估儿童行为问题和情绪障碍。
            \item \textbf{个性测验}：如艾森克人格问卷（EPQ）、儿童气质问卷等。
            \item \textbf{神经心理测验}：评估儿童注意、记忆、执行功能等特定认知领域的功能。
        \end{itemize}

    \item \imp{流行病学调查方法}——运用流行病学研究设计，探讨儿童少年健康状况及其影响因素：
        \begin{itemize}
            \item \textbf{横断面调查（cross-sectional survey）}：在特定时间点收集儿童少年群体的健康资料，描述生长发育水平、疾病分布状况、健康行为流行特征等，是学校卫生监测最常用的方法。
            \item \textbf{追踪调查（longitudinal study）}：对同一群体进行长期、连续的动态观察，揭示生长发育的轨迹和变化规律，分析早期因素对远期健康结局的影响，是研究生长发育规律和因果关联的重要方法。
            \item 此外还包括病例对照研究（case-control study）和队列研究（cohort study）等分析性流行病学方法。
        \end{itemize}

    \item \imp{统计分析方法}——对收集的数据进行科学的统计处理，主要包括：
        \begin{itemize}
            \item \textbf{描述性统计}：均值、标准差、百分位数、生长曲线拟合等。
            \item \textbf{推断性统计}：t检验、方差分析、卡方检验、相关与回归分析等。
            \item \textbf{多因素分析}：多元线性回归、Logistic回归、Cox比例风险模型等。
            \item \textbf{纵向数据分析}：重复测量方差分析、混合效应模型、生长曲线模型等，专门处理追踪数据中的时间效应和个体间差异。
            \item \textbf{生长评价方法}：Z-score法、百分位数法、最小二乘均值法等标准化评价技术。
        \end{itemize}
\end{enumerate}

\textbf{方法学的综合运用：}在实际研究中，上述方法常综合应用，形成多维度、多方法的综合评价体系。如\"人体测量+生理功能检测+心理行为评定\"相结合的生长发育综合评价，\"流行病学调查+统计分析\"相结合的健康状况和影响因素研究等。
\end{defbox}

\subsection{三级预防理论}

\begin{keybox}
\textbf{三级预防理论：}

三级预防理论是学科发展的重要理论构架。三级预防是儿童少年健康促进的首要和有效手段，是现代医学为人们提供的健康保障。

\begin{enumerate}[leftmargin=*]
    \item \imp{一级预防（primary prevention）/ 病因预防}：疾病尚未发生时针对致病因素或危险因素采取措施，是预防和消灭疾病的根本措施。是最积极最有效的预防措施。

    \item \imp{二级预防（secondary prevention）/ "三早"预防}：早发现、早诊断、早治疗，是防止或减缓疾病发展采取的措施。

    \item \imp{三级预防（tertiary prevention）/ 临床预防}：防治伤残和促进功能恢复，提高生命质量，延长寿命，降低病死率，主要是对症治疗和康复治疗措施。
\end{enumerate}
\end{keybox}

\newpage
